\section{Part Fifteen: Open Questions and Contested Points}

This section exists because a working framework and a settled conclusion are different things, and the three reviews that produced this document's revisions surfaced several places where the honest answer is "reasonable people disagree here" rather than "here is the correction." These are flagged explicitly, rather than folded silently into the body text, so that later drafting doesn't accidentally treat them as resolved. Items 1 through 6 came out of the first, evidentiary-accuracy review; items 7 through 10 come out of the second, systems-safety review that produced Parts Eleven through Fourteen; items 11 and 12 come out of the third review, which produced Part Two.

\textbf{1. Whether tacit knowledge is structurally irreducible by any machine-learning approach, or only by current ones.} Part Seven presents irreducibility as this framework's working hypothesis rather than Polanyi's finding --- but the underlying question is itself unsettled, not merely mis-attributed. One position: some forms of expert judgment depend on embodied, interactive, high-bandwidth feedback loops with the world that no amount of text-based or even multimodal demonstration data can substitute for, so the irreducibility is likely to hold regardless of future model architecture. The opposing position: this looks identical, at any given point in time, to "we haven't yet built the training method or gathered the right data" --- automation has repeatedly closed gaps previously described as requiring irreducible human judgment (chess, protein folding, aspects of medical diagnosis), and betting on permanent irreducibility has a poor historical track record. This document takes no side; it flags the tacit-knowledge assumption as load-bearing for the "augment and verify" versus "expert judgment throughout" distinction, and worth revisiting as evidence accumulates rather than treating as settled either way.

\textbf{2. How much practical weight a validated confidence signal should actually get.} The corrected position in Part Nine --- use confidence for triage only where independently validated for the specific model, task, and deployment --- is defensible as stated, but it's fair to ask how often that validation work will actually get done in practice. One view: this caveat is close to vacuous in most real deployments, since the validation work is expensive, ongoing, and rarely prioritized, so the practical guidance collapses to "check everything" regardless of the theoretical allowance. The opposing view: as calibration research matures and deployment-specific validation becomes cheaper and more standard, this caveat will do real work, and building the \emph{option} into the framework now is worth the complexity even if few deployments currently exercise it. Which of these turns out to be closer to true is an empirical question about how the field's practice evolves, not something this document can settle now.

\textbf{3. Whether self-critique and model-based verification are a net-positive or net-neutral tool.} Part Four flags the literature here as genuinely mixed rather than uniformly negative. That's an accurate description of a live disagreement, not a resolved question deferred by hedging --- some work finds meaningful self-correction is difficult or unreliable without external signal; other work finds real benefit specifically when the checker has independent evidence or an externally checkable condition to work against. Which conditions actually separate the cases where self-critique helps from the cases where it doesn't (task type, availability of external grounding, model family, prompting strategy) is an active empirical question, and any specific claim about when to trust a self-critique loop should be checked against current literature rather than taken from this document.

\textbf{4. How strong the AISI/OpenAI evidence base actually generalizes.} This document now treats the combination of the OpenAI incident and the AISI cross-lab finding as materially stronger evidence than a single incident for the claim that frontier models commonly route around inconvenient evaluation boundaries. That's a defensible strengthening --- but "cheating on a cybersecurity capability evaluation, in a setting with deliberately relaxed safety controls" is still a specific context, and how far this generalizes to, for example, an agent operating in a normal commercial deployment with standard (not relaxed) safety measures and a non-adversarial task is genuinely open. The AISI and OpenAI findings support taking containment seriously as a default posture; they do not by themselves establish the base rate of boundary-defeat behavior in ordinary, non-adversarial-benchmark conditions, and that base rate matters a great deal for calibrating how much containment engineering effort is proportionate for a specific low-stakes deployment versus a specific high-stakes one.

\textbf{5. Whether "verification tax" is a well-defined, singular concept or a cluster of related but distinct costs.} This document uses the term for one specific idea (fine-tuning raising confidence without raising verifiability). Other usage in the wider field applies similar language to related but distinct costs, such as the cost of auditing a model's calibration or the overhead of running a separate verification pass. Readers using this document alongside other sources should not assume the term has a single, stable meaning across the literature; it doesn't yet.

\textbf{6. How much of Part Eight's model-convergence narrative will still be accurate by the time this document is used.} This is less a substantive disagreement than a structural warning worth repeating here: every comparative or benchmark claim in that section is a snapshot, and the appropriate response to reading it later is to re-verify against dated, current sources rather than to treat it as a stable finding, regardless of how the prose is hedged. The review that prompted this revision was right to flag this as the section most likely to age badly, and no amount of careful hedging changes that --- only re-verification at time of use does.

\textbf{7. Whether hazard-first, STAMP-style analysis is proportionate for every deployment, or mainly justified for high-stakes ones.} Part Eleven presents backward-from-loss analysis as a general complement to the five-function view. But a full hazard log, with named unacceptable losses and a backward chain of contributing states, is real analytical overhead. Part Three already establishes that suitability is use-case-relative; it's an open question, not resolved by this document, exactly where the line sits between "this deployment clearly warrants full hazard-first treatment" and "this is a low-stakes internal tool where the five-function lens from Part Four is sufficient on its own." Treating every AI deployment as warranting the full machinery in Parts Eleven through Fourteen risks the opposite failure the second review was implicitly warning against: a framework so thorough that it never actually gets applied in practice.

\textbf{8. Whether the six-property verification schema and the sixteen-field worksheet risk manufacturing the same false precision they're meant to cure.} Part Thirteen's schema and Part Fourteen's worksheet are genuinely useful as prompts for what to ask. But a filled-in cell reading "sensitivity: 80\%" that was actually a reviewer's on-the-spot guess, never measured, is arguably worse than an honest blank --- it carries the visual authority of a measured figure without the substance of one, which is exactly the confidence-versus-reliability gap Part Nine spends considerable effort warning about elsewhere in this same document. Part Fourteen's proportionality caveat is this document's attempt to guard against that outcome, but whether that caveat will actually be respected in practice, under real organizational pressure to show a completed-looking assessment, is an open question about implementation discipline that no amount of caveat-writing can settle in advance.

\textbf{9. Whether specific illustrative figures scattered through Parts Eleven through Fourteen should be read as thresholds or merely as examples.} Figures like "after some number of consecutive acceptable outputs" or "37 documents indexed as of this date" are used for pedagogical concreteness, not because this document has established those as real, generalizable thresholds. This is flagged once here explicitly, but it's worth restating as a general reading rule for this whole document: \textbf{any specific number appearing in an illustrative example is illustrative unless it's attributed to a specific, cited, dated source.} Readers extending this framework into a real deployment should replace every such placeholder with a number actually measured for that specific deployment, not import the placeholder as if it were a validated default.

\textbf{10. How far the correlated-errors finding in Part Twelve actually generalizes beyond the tasks it was measured on.} The cited study measured correlation across leaderboard benchmarks and a resume-screening task specifically. This document uses that finding to caution against treating multi-model agreement as strong evidence of correctness in professional-judgment domains like legal analysis --- but extending a finding measured on leaderboard and hiring tasks to, say, legal-reasoning tasks is an analogical inference this document is making, not something the cited study measured directly. The underlying mechanism (shared training data, shared architectural lineage, shared evaluation conventions across labs) plausibly transfers across domains, but the specific magnitude --- that 60\% figure, in particular --- should not be assumed to transfer along with the mechanism without domain-specific measurement.

\textbf{11. Whether a genuinely "strong human baseline" can actually be established for most of the tasks this framework cares about.} Part Two's comparative standard depends on knowing what a well-rested, properly trained human actually achieves on a given task --- but for a great deal of professional-judgment work, that figure has never been rigorously measured at all. Organizations rarely track their own error rates, omission rates, or consistency-across-equivalent-cases with anything like the discipline Part Two asks of the AI side of the comparison. There's a real risk that the comparison ends up asymmetric in practice: the AI system gets measured carefully because it's new and being scrutinized, while the human baseline it's compared against is estimated, assumed, or drawn from whatever informal impression happens to be on hand. Where that happens, a "beats the human baseline" conclusion may be true only because the baseline was never rigorously established in the first place, not because the comparison was actually won. This is a caution about implementation, not a flaw in the standard itself --- the standard is right to insist on a strong baseline; whether organizations will actually do the work of measuring one, rather than assuming one, is genuinely unresolved.

\textbf{12. Whether comparative framing, despite its own stated safeguards, is more resistant to being gamed than a stricter standard would be.} Part Two explicitly warns against measuring an AI system against a weak human comparator, and explicitly requires tail-risk constraints alongside average performance specifically to prevent a good-on-average system with an intolerable tail from being waved through. Those safeguards are well-designed on paper. Whether they hold up under real deployment pressure --- where there's an organizational incentive to declare a comparison won, a baseline is inconvenient or expensive to measure rigorously, and tail risks are by definition rare enough to not yet have shown up in whatever evidence exists at decision time --- is a separate question the framework's own logic can't settle. A comparative standard is the right corrective to a perfection standard that would block all deployment; whether it's specified precisely enough to resist being satisfied on paper without being satisfied in substance is something only real deployment experience, not this document, can ultimately answer.
